\chapter*{Fiche synthétique du stage}
\addcontentsline{toc}{chapter}{Fiche synthétique du stage}
Ce rapport présente l’introduction générale et les quatre premiers chapitres du travail de stage : le cadrage, le benchmark, l’analyse des besoins et l’architecture fonctionnelle cible. La modélisation UML détaillée et l’architecture applicative logique constituent les prolongements prévus de cette démarche de conception.

\begin{tabularx}{\linewidth}{@{}P{37mm}X@{}}
\toprule
\textbf{Sujet} & Conception d’une architecture intelligente pour transformer un WMS traditionnel en Smart WMS\\[8pt]
\textbf{Nature du travail} & Conception, benchmark, analyse fonctionnelle, préparation de la modélisation UML et d’une architecture applicative logique\\[8pt]
\textbf{Approche} & SCAN : observer et comparer ; PLAN : structurer et concevoir ; ACT : formaliser et vérifier les scénarios\\[8pt]
\textbf{Solutions étudiées} & Manhattan Active WMS, SAP EWM et Oracle WMS Cloud\\[8pt]
\textbf{Orientation} & Architecture générique, modulaire, adaptable et configurable\\[8pt]
\textbf{Résultat attendu} & Architecture de référence adaptable aux contraintes métier, aux données et au niveau d’infrastructure de différents entrepôts\\
\bottomrule
\end{tabularx}
\clearpage
